\subsection{Implementasi Segmentasi Data RME}
\label{subsec:implementasi-segmentasi-data-rme}

% Implementasi segmentasi data RME pada subbab ini merupakan realisasi dari \texttt{R-TA-02} melalui komponen \texttt{I-TA-01}. Bagian metadata \textit{on-chain} yang mendukung segmentasi juga terkait dengan \texttt{I-TA-04}.

Segmentasi mengikuti rancangan pada Subbab \ref{subsec:rancangan-segmentasi-data-rme} dan diwujudkan pada \textit{crate} \texttt{decmed-rme-segment}. \textit{Crate} ini mendefinisikan \texttt{DatasetCategory} dan \texttt{FunctionCategory} sebagai atribut objek untuk penegakan \textit{policy}. Rincian modul ada pada Lampiran \ref{lampiran:implementasi-crate-decmed-rme-segment}. Implementasi enum ditunjukkan pada Kode \ref{lst:implementasi-dataset-category-segmentasi-rme} dan \ref{lst:implementasi-function-category-segmentasi-rme}.

\begin{lstlisting}[language=Rust, caption=Implementasi \texttt{DatasetCategory} pada \textit{crate} \texttt{decmed-rme-segment}, label={lst:implementasi-dataset-category-segmentasi-rme}]
	pub enum DatasetCategory {
		RAWAT_JALAN,
		RAWAT_INAP,
		LABORATORIUM,
		APOTEK,
	}
\end{lstlisting}

\begin{lstlisting}[language=Rust, caption=Implementasi \texttt{FunctionCategory} pada \textit{crate} \texttt{decmed-rme-segment}, label={lst:implementasi-function-category-segmentasi-rme}]
	pub enum FunctionCategory {
		// ADMINISTRATIVE
		ADMINISTRATIVE_GENERAL,
	
		// RAWAT_JALAN & RAWAT_INAP
		ANAMNESIS,
		PEMERIKSAAN_FISIK,
		PEMERIKSAAN_PSIKOLOGIS,
		RIWAYAT_PENGGUNAAN_OBAT,
		RENCANA_RAWAT,
		INSTRUKSI_MEDIK_DAN_KEPERAWATAN,
		PEMERIKSAAN_PENUNJANG,
		DIAGNOSIS,
		INFORMED_CONSENT,
		TERAPI,
	
		// RAWAT_INAP
		PERENCANAAN_PEMULANGAN,
	
		// LABORATORIUM
		LABORATORIUM,
	
		// APOTEK
		PERESEPAN,
		DISPENSING,
	}
\end{lstlisting}

\subsubsection{Implementasi Segmentasi Metadata RME pada Penyimpanan \textit{On-Chain}}
\label{subsec:implementasi-segmentasi-metadata-rme-onchain}

\texttt{RmeSegmentMetadata} merepresentasikan metadata segmen sebelum diserialisasi menjadi JSON base64 dan disimpan pada \textit{field} \texttt{metadata} milik \texttt{PatientMedicalMetadata} (Kode \ref{lst:decmed-patient-medical-metadata}). Struktur mengikuti skema pada Tabel~\ref{tab:onchain-schema} di Lampiran~\ref{lampiran:skema-penyimpanan-rme}; implementasinya pada Kode~\ref{lst:implementasi-segmentasi-metadata-rme}.

\begin{lstlisting}[language=Rust, caption=Implementasi struktur \texttt{RmeSegmentMetadata} pada \textit{crate} \texttt{decmed-rme-segment}, label={lst:implementasi-segmentasi-metadata-rme}]
pub struct RmeSegmentMetadata {
	pub segment_id: String,
	pub related_rme_id: String,
	pub patient_address: String,
	pub author_address: String,
	pub hospital_cid: String,
	pub dataset_category: DatasetCategory,
	pub function_category: FunctionCategory,
	pub ipfs_cid: String,
	pub integrity_hash: String,
	pub capsule: String,
	pub enc_key_and_nonce: String,
	pub created_at: String,
}
\end{lstlisting}

Sebelum disimpan ke IOTA, metadata divalidasi: setiap \textit{field} wajib terisi, kombinasi \texttt{dataset\_category} dan \texttt{function\_category} sah, serta \texttt{integrity\_hash} sesuai data terenkripsi.

\subsubsection{Implementasi Segmentasi Data RME pada Penyimpanan \textit{Off-Chain}}
\label{subsec:implementasi-segmentasi-data-rme-offchain}

\texttt{RmeSegmentData} merepresentasikan segmen sebelum enkripsi di sisi \textit{client}, sesuai skema pada Tabel~\ref{tab:offchain-schema} di Lampiran~\ref{lampiran:skema-penyimpanan-rme}. Implementasinya pada Kode~\ref{lst:implementasi-segmentasi-data-rme}. Pada implementasi ini, \texttt{payload} bertipe \texttt{Value} dan berisi data berbentuk teks; \texttt{payload\_hash} adalah hash integritas dari konten \texttt{payload}.

\begin{lstlisting}[language=Rust, caption=Implementasi struktur \texttt{RmeSegmentData} pada \textit{crate} \texttt{decmed-rme-segment}, label={lst:implementasi-segmentasi-data-rme}]
pub struct RmeSegmentData {
    pub segment_id: String,
    pub related_rme_id: String,
    pub dataset_category: DatasetCategory,
    pub function_category: FunctionCategory,
    pub patient_address: String,
    pub service_date: String,
    pub author_address: String,
    pub payload: Value,
    pub payload_hash: String,
}
\end{lstlisting}

Saat penyimpanan, \textit{Hospital Client} membentuk objek per segmen, mengenkripsinya, mengunggah ke IPFS, lalu menyimpan CID bersama metadata pada IOTA sebagai JSON base64 di \textit{field} \texttt{metadata} milik \texttt{PatientMedicalMetadata}. Pendekatan ini menjaga kompatibilitas dengan daftar metadata \textit{flat} DecMed, sambil memungkinkan PRE membaca kategori segmentasi dari isi metadata.

% \texttt{payload\_hash} dihitung dari isi \texttt{payload} sebelum enkripsi, sedangkan \texttt{integrity\_hash} dihitung dari data segmen terenkripsi yang disimpan di IPFS. Dengan demikian, \texttt{payload\_hash} digunakan untuk integritas isi segmen sebelum enkripsi, sedangkan \texttt{integrity\_hash} digunakan untuk memverifikasi kesesuaian \textit{ciphertext} IPFS dengan metadata pada IOTA.

% Pada implementasi ini, \texttt{RmeSegmentData} digunakan sebagai representasi segmen sebelum enkripsi pada sisi \textit{client}. Setelah \texttt{payload} dienkripsi, segmen direpresentasikan sebagai \texttt{ClientEncryptedRmeSegment} yang masih memuat \textit{ciphertext}. Struktur tersebut kemudian dikonversi menjadi \texttt{RmeSegmentMetadata} dengan menghilangkan isi \textit{ciphertext} dan mempertahankan metadata yang diperlukan untuk pencarian, validasi integritas, dan proses PRE. Metadata ini kemudian diserialisasi sebagai JSON base64 sebelum disimpan pada field \texttt{metadata} di IOTA.

% Implementasi ini membuat proses otorisasi dapat dilakukan pada tingkat segmen. Sebagai contoh, token Macaroon yang hanya memuat izin baca untuk \texttt{RAWAT\_JALAN}, fungsi \texttt{PEMERIKSAAN\_PENUNJANG}, dan \texttt{related\_rme\_id} tertentu hanya dapat digunakan untuk meminta segmen yang berada pada episode dan kategori tersebut. Segmen lainnya seperti \texttt{DIAGNOSIS}, \texttt{TERAPI}, atau \texttt{PERESEPAN} tidak akan dikirimkan ke \textit{client} oleh \textit{Server} PRE karena metadata segmennya tidak termasuk dalam cakupan efektif Macaroon.

% Dengan pendekatan ini, IOTA berperan sebagai penyimpanan metadata dan sumber rujukan status segmen, sedangkan penegakan aturan akses berdasarkan kategori dilakukan oleh \textit{Server} PRE pada saat permintaan baca atau tulis diproses.

% Segmentasi juga digunakan pada proses pembaruan data. Ketika aktor menulis data baru, \textit{Hospital Client} hanya membuat atau memperbarui segmen yang berada dalam cakupan izin tulis aktor tersebut. Dengan demikian, petugas laboratorium dapat menambahkan hasil pemeriksaan ke segmen data laboratorium tanpa memperoleh akses tulis terhadap segmen diagnosis atau terapi. Pembaruan segmen direpresentasikan dengan perubahan metadata, misalnya perubahan \texttt{ipfs\_cid} dan pengisian \texttt{updated\_at}, tanpa memberikan hak tulis terhadap segmen lain di luar cakupan token.
